\documentclass[11pt]{article}
\usepackage[margin=1in]{geometry}
\usepackage{amsmath,amssymb,booktabs}
\usepackage[T1]{fontenc}
\usepackage[utf8]{inputenc}

\title{Welfare Normalization in the Tidal-Lane Spatial Equilibrium}
\author{}
\date{July 1, 2026}

\begin{document}
\maketitle

\section{Object}

The primitive welfare index in the model is the Frechet denominator over residence, workplace, and routes. After aggregating routes into the bilateral commuting-cost index, the object is
\begin{equation}
\label{eq:primitive_welfare}
    W^\theta
    =
    \sum_i \sum_j
    \tau_{ij}^{-\theta} u_i^\theta w_j^\theta ,
\end{equation}
where
\[
    u_i^\theta = \bar u_i^\theta (L_i^R)^{\beta\theta},
    \qquad
    w_j^\theta = \bar A_j^\theta (L_j^F)^{\alpha\theta}.
\]
Equation \eqref{eq:primitive_welfare} is invariant to any normalization of the resident and firm market-access indices. It is therefore the correct object for comparing welfare across counterfactual equilibria.

\section{Market-Access Representation}

The solver also computes resident and firm market-access indices from
\begin{align}
\label{eq:market_access_pi}
    \Pi_i^{-\theta}
    &=
    \sum_j \tau_{ij}^{-\theta} L_j^F P_j^\theta, \\
\label{eq:market_access_p}
    P_j^{-\theta}
    &=
    \sum_i \tau_{ij}^{-\theta} L_i^R \Pi_i^\theta .
\end{align}
This system does not determine the absolute scale of $(\Pi,P)$. If $(\Pi,P)$ solves \eqref{eq:market_access_pi}--\eqref{eq:market_access_p}, then for any $c>0$,
\[
    \tilde \Pi_i = c\Pi_i,
    \qquad
    \tilde P_j = P_j/c
\]
is also a solution. The transformation leaves link traffic invariant because
\[
    \tilde P_k^{-\theta}\tilde \Pi_l^{-\theta}
    =
    c^\theta P_k^{-\theta}
    c^{-\theta} \Pi_l^{-\theta}
    =
    P_k^{-\theta}\Pi_l^{-\theta}.
\]
It follows that edge times, bilateral costs, link traffic, residence shares, workplace shares, and OD-weighted commuting time are unaffected by this scale choice.

\section{Why the Previous Welfare Number Drifted}

The previous implementation evaluated
\begin{equation}
\label{eq:pi_welfare}
    W_\Pi^\theta = \sum_i u_i^\theta \Pi_i^{-\theta}.
\end{equation}
This expression is invariant only after the market-access gauge is fixed consistently across scenarios. Under $\Pi_i \mapsto c\Pi_i$,
\[
    W_\Pi^\theta \mapsto c^{-\theta}W_\Pi^\theta.
\]
The fixed-point routine initialized $(\Pi,P)$ at $(1,1)$ in each scenario. That procedure selects a deterministic scale, but the selected scale is scenario-specific and has no economic content. A cross-scenario comparison of \eqref{eq:pi_welfare} can therefore combine the primitive welfare change with a pure normalization change.

The market-access representation is equivalent to the primitive welfare index only under a gauge that maps firm market access into wages:
\begin{equation}
\label{eq:gauge_condition}
    L_j^F P_j^\theta = w_j^\theta
    \quad \text{for all } j.
\end{equation}
When \eqref{eq:gauge_condition} holds,
\[
    \Pi_i^{-\theta}
    =
    \sum_j \tau_{ij}^{-\theta} w_j^\theta ,
\]
and substituting this expression into \eqref{eq:pi_welfare} recovers \eqref{eq:primitive_welfare}. Without \eqref{eq:gauge_condition}, \eqref{eq:pi_welfare} has a normalization component.

\section{Correct Implementation}

The implementation should compute welfare directly from
\begin{equation}
\label{eq:implemented_welfare}
    W
    =
    \left[
    (u^\theta)^\top
    (I-A)^{-1}
    w^\theta
    \right]^{1/\theta},
\end{equation}
with
\[
    A_{kl} = (t_{kl})^{-\theta},
    \qquad
    (I-A)^{-1}_{ij} = \tau_{ij}^{-\theta}.
\]
This formula uses objects already solved in each equilibrium and does not require $\Pi$ or $P$. A gauge-corrected resident market-access map can still be reported, but it should be computed as
\begin{equation}
\label{eq:gauge_free_pi}
    \Pi_{i,w}^{-\theta}
    =
    \sum_j \tau_{ij}^{-\theta} w_j^\theta .
\end{equation}
Maps based on the raw $\Pi$ from \eqref{eq:market_access_pi} should be interpreted only after removing the scenario-level scale component.

\section{Numerical Implication for the Main Tidal-Lane Run}

For the gate5 BPR-selected top-100 tidal-lane counterfactual with $(\alpha,\beta)=(0.05,-0.10)$, the old market-access welfare number and the primitive welfare number differ in the long run:
\begin{center}
\begin{tabular}{lcc}
\toprule
Welfare formula & Short run & Long run \\
\midrule
$\sum_i u_i^\theta\Pi_i^{-\theta}$ with raw solver gauge
    & $+0.94\%$ & $-0.85\%$ \\
$\sum_{ij}\tau_{ij}^{-\theta}u_i^\theta w_j^\theta$
    & $+0.87\%$ & $+0.32\%$ \\
\bottomrule
\end{tabular}
\end{center}
The OD-weighted commuting-cost decomposition remains a separate object. In the same run, the average bilateral commuting cost falls in the short run and is nearly unchanged in the long run: the long-run total is about $+0.065$ seconds per commuter, equal to the sum of a negative direct effect and positive general-equilibrium and sorting wedges.

\section{Reporting Rule}

All welfare tables and slide bullets should report \eqref{eq:implemented_welfare}. Raw $\Pi$ maps should be replaced by either the gauge-free market-access object in \eqref{eq:gauge_free_pi} or a demeaned change in $\log \Pi$ that removes the scenario-level normalization. Population sorting, workplace sorting, link traffic, travel time, and person-minute results do not require this correction.

\end{document}
